\section{Related Work}
\label{sec:related}

\para{Causal interventions in mechanistic interpretability.}
The dominant approach to understanding neural network internals uses causal interventions on model activations. \citet{heimersheim2024interpret} provide a systematic framework connecting intervention direction to causal claims: noising (replacing clean activations with corrupted ones) tests necessity, while denoising (the reverse) tests sufficiency. They show these are not symmetric---AND-structured circuits are fully discoverable via noising, while OR-structured circuits require denoising. \citet{zhang2023best} compare corruption strategies and show that symmetric token replacement (\str) avoids the off-distribution artifacts of Gaussian noise. Our work builds on both: we use \str corruption and test both intervention directions, but uniquely evaluate whether identified components are specific to one task or affect many.

\para{Circuit discovery and evaluation.}
The Indirect Object Identification circuit \citep{wang2022interpretability} is a landmark result, identifying attention heads that move name information to predict indirect objects. Subsequent work has extended circuit-level analysis to subject-verb agreement, greater-than reasoning, and other tasks. The Mechanistic Interpretability Benchmark \citep[\textsc{MIB};][]{mueller2025mib} formalizes circuit evaluation with two metrics: Circuit Proportion Recovered (\textsc{CPR}, measuring sufficiency) and Circuit Miss Discovery (\textsc{CMD}, measuring necessity). However, neither \textsc{MIB} nor the original circuit papers measure whether identified components are specific to their target task. We directly address this gap.

\para{Selectivity in interpretability.}
\citet{mueller2024quest} identify four evaluation criteria for causal mediators: sparsity, generality, selectivity, and faithfulness. They note that selectivity---explaining only the target behavior and as little else as possible---is ``not often explicitly discussed nor measured.'' \citet{arora2024causalgym} introduce a selectivity metric for \textsc{CausalGym}, computing the difference between a method's causal efficacy on real tasks versus control tasks. They find that after selectivity adjustment, linear probing outperforms DAS at larger model scales. Unlike \textsc{CausalGym}, which evaluates 1D linear features on linguistic tasks, we evaluate full model components (attention heads and MLP layers) across tasks from different domains.

\para{Feature-level interpretability.}
An alternative to component-level analysis operates at the feature level. \citet{marks2024sparse} discover circuits using SAE features and demonstrate selectivity in practice: ablating task-irrelevant features removes gender bias while preserving profession classification. However, \citet{mueller2025mib} find that SAE features do not improve over standard neurons for causal variable localization, questioning whether feature decomposition reliably improves selectivity. \citet{sharkey2025open} identify fundamental limitations of sparse dictionary learning, including feature splitting and absorption. Our work complements this line by showing that even at the component level, selectivity failures are systematic and quantifiable.

\para{Causal abstraction.}
\citet{geiger2023causal} provide a formal framework grounding mechanistic interpretability in causal abstraction theory, defining interchange intervention accuracy (\textsc{IIA}) as a measure of counterfactual faithfulness. \citet{bereska2024review} survey the field through an AND/OR logic lens, connecting circuit structure to the relative informativeness of necessity and sufficiency tests. Our empirical results align with these theoretical predictions: \ioi shows the most divergence between necessity and sufficiency ($\rho = 0.760$), consistent with its known OR-circuit structure with redundant name mover heads.